% =============================================================================
% header_compiladores.tex — Header acadêmico | Projeto de Compiladores
% Base: header_4.tex (emoji, mono, zebra, setrightheader/setrightsub)
%
% Uso:
%   pandoc ARQUIVO.md -o ARQUIVO.pdf \
%     --pdf-engine=lualatex \
%     --include-in-header=header_compiladores.tex
%
% No .md (opcional — lado direito do cabeçalho):
%   \setrightheader{Análise Léxica}     % título à direita
%   \setrightsub{Fase 1 · 2026.2}      % subtítulo à direita
% Sem \setrightheader / \setrightsub: o lado direito fica em branco.
% =============================================================================

\usepackage{fontspec}
\usepackage{geometry}
\usepackage{graphicx}
\usepackage{fancyhdr}
\usepackage[table]{xcolor}
\usepackage{colortbl}
\usepackage{array}
\usepackage{booktabs}
\usepackage{setspace}
\usepackage{titlesec}
\usepackage{longtable}
\usepackage{etoolbox}

% -----------------------------------------------------------------------------
% Emoji — fallback para Noto Color Emoji
% -----------------------------------------------------------------------------
\directlua{
luaotfload.add_fallback("emojifallback",{
  "Noto Color Emoji:mode=harf;"
})
}

% -----------------------------------------------------------------------------
% Cores
% -----------------------------------------------------------------------------
\definecolor{corprincipal}{HTML}{124860}
\definecolor{corsecundaria}{HTML}{84AFBF}
\definecolor{cinzaclaro}{HTML}{F2F4F5}
\definecolor{cinzamedio}{HTML}{E6EBED}
\definecolor{texto}{HTML}{1A1A1A}
\definecolor{textoleve}{HTML}{5A6A70}
\definecolor{azulclaro}{HTML}{B1D7F2}

% -----------------------------------------------------------------------------
% Página
% -----------------------------------------------------------------------------
\geometry{
  a4paper,
  top=2.8cm,
  bottom=2.4cm,
  left=2cm,
  right=2cm,
  headheight=1.6cm,
  headsep=0.6cm,
  footskip=1.2cm
}

% -----------------------------------------------------------------------------
% Tipografia — Lufga + emoji fallback + JetBrains Mono
% -----------------------------------------------------------------------------
\setmainfont{Lufga}[
  UprightFont=* Regular,
  BoldFont=* Bold,
  ItalicFont=* Italic,
  BoldItalicFont=* Bold Italic,
  Ligatures=TeX,
  RawFeature={-liga;fallback=emojifallback}
]

\setsansfont{Lufga}[
  UprightFont=* Regular,
  BoldFont=* Bold,
  ItalicFont=* Italic,
  BoldItalicFont=* Bold Italic,
  Ligatures=TeX,
  RawFeature={-liga;fallback=emojifallback}
]

% Scale=1 evita o recuo típico do mono (fica ~10–15% menor que o corpo)
\setmonofont{JetBrains Mono}[Scale=1.0]

% -----------------------------------------------------------------------------
% Blocos de código (Pandoc: Highlighting + Shaded + verbatim + \texttt)
% Shaded só é criado pelo template do Pandoc — ajustar em \AtBeginDocument.
% -----------------------------------------------------------------------------
\usepackage{fancyvrb}
\usepackage{framed}
\definecolor{shadecolor}{HTML}{F2F4F5}

\makeatletter
\renewcommand{\verbatim@font}{\normalfont\normalsize\ttfamily}
\makeatother
\let\oldtexttt\texttt
\renewcommand{\texttt}[1]{{\normalsize\ttfamily #1}}

% Não mexe em Shaded aqui (o Pandoc define depois; renew no preâmbulo quebra).
\AtBeginDocument{%
  \fvset{
    fontsize=\normalsize,
    baselinestretch=1.15,
    formatcom=\color{texto}
  }%
  \DefineVerbatimEnvironment{Highlighting}{Verbatim}{%
    fontsize=\normalsize,
    commandchars=\\\{\},
    baselinestretch=1.15,
    formatcom=\color{texto}
  }%
}

% -----------------------------------------------------------------------------
% Tabelas
% Pandoc longtable: corpo zebra
% tabular: zebra a partir da 2ª linha
% -----------------------------------------------------------------------------
\renewcommand{\arraystretch}{1.3}
\setlength{\tabcolsep}{8pt}
\setlength{\aboverulesep}{0pt}
\setlength{\belowrulesep}{0pt}
\setlength{\extrarowheight}{1pt}
\arrayrulecolor{corsecundaria}

\makeatletter
\def\toprule\noalign#1{%
  \color{corprincipal}\bfseries
}
\makeatother

\AtBeginEnvironment{longtable}{%
  \rowcolors{2}{white}{cinzaclaro}%
  \arrayrulecolor{corsecundaria}%
}

\AtBeginEnvironment{tabular}{%
  \rowcolors{2}{azulclaro}{white}%
  \arrayrulecolor{corsecundaria}%
}

% -----------------------------------------------------------------------------
% Lado direito do cabeçalho
%   \setrightheader{...}  → título
%   \setrightsub{...}     → subtítulo
% -----------------------------------------------------------------------------
\newcommand{\rightheadertext}{}
\newcommand{\rightsubtext}{}
\newcommand{\setrightheader}[1]{\renewcommand{\rightheadertext}{#1}}
\newcommand{\setrightsub}[1]{\renewcommand{\rightsubtext}{#1}}

% -----------------------------------------------------------------------------
% Lado esquerdo — texto do projeto (sem logo)
% -----------------------------------------------------------------------------
\newcommand{\leftheadertext}{Projeto de Compiladores}
\newcommand{\leftsubtext}{Documento acadêmico}
\newcommand{\setleftheader}[1]{\renewcommand{\leftheadertext}{#1}}
\newcommand{\setleftsub}[1]{\renewcommand{\leftsubtext}{#1}}

\newcommand{\leftheadblock}{%
  \begin{minipage}[b]{8.2cm}
    {\fontsize{11}{13}\selectfont\color{corprincipal}\textbf{\leftheadertext}}%
    \ifdefempty{\leftsubtext}{}{%
      \\[-0.12em]
      {\fontsize{7}{9}\selectfont\color{textoleve}\leftsubtext}%
    }%
  \end{minipage}%
}

\newcommand{\rightheadblock}{%
  \begin{minipage}[b]{5.6cm}
    \raggedleft
    \ifdefempty{\rightheadertext}{}{%
      {\fontsize{10}{12}\selectfont\color{corprincipal}\textbf{\rightheadertext}}%
    }%
    \ifdefempty{\rightsubtext}{}{%
      \\[-0.15em]
      {\fontsize{7}{9}\selectfont\color{textoleve}\rightsubtext}%
    }%
  \end{minipage}%
}

% -----------------------------------------------------------------------------
% Header e Footer
% -----------------------------------------------------------------------------
\pagestyle{fancy}
\fancyhf{}

\fancyhead[L]{\leftheadblock}
\fancyhead[R]{\rightheadblock}

\renewcommand{\headrulewidth}{1.2pt}
\renewcommand{\headrule}{%
  \hbox to\headwidth{%
    \color{corprincipal}\leaders\hrule height \headrulewidth\hfill
  }%
}

\fancyfoot[L]{%
  {\fontsize{7.5}{9}\selectfont\color{textoleve}%
    Projeto de Compiladores%
    \ifdefempty{\rightheadertext}{}{ · \rightheadertext}}%
}
\fancyfoot[C]{}
\fancyfoot[R]{%
  {\fontsize{8}{10}\selectfont\color{corprincipal}\textbf{\thepage}}%
}

\renewcommand{\footrulewidth}{0.6pt}
\renewcommand{\footrule}{%
  \hbox to\headwidth{%
    \color{corsecundaria}\leaders\hrule height \footrulewidth\hfill
  }%
}

\fancypagestyle{plain}{%
  \fancyhf{}
  \fancyhead[L]{\leftheadblock}
  \fancyhead[R]{\rightheadblock}
  \fancyfoot[L]{%
    {\fontsize{7.5}{9}\selectfont\color{textoleve}%
      Projeto de Compiladores%
      \ifdefempty{\rightheadertext}{}{ · \rightheadertext}}%
  }
  \fancyfoot[R]{%
    {\fontsize{8}{10}\selectfont\color{corprincipal}\textbf{\thepage}}%
  }
  \renewcommand{\headrulewidth}{1.2pt}
  \renewcommand{\footrulewidth}{0.6pt}
}

% -----------------------------------------------------------------------------
% Títulos de seção
% -----------------------------------------------------------------------------
\titleformat{\section}
  {\color{corprincipal}\Large\bfseries}
  {}{0em}{}
  [\vspace{0.15em}{\color{corsecundaria}\titlerule[1.1pt]}]

\titleformat{\subsection}
  {\color{corprincipal}\large\bfseries}
  {}{0em}{}

\titleformat{\subsubsection}
  {\color{corprincipal}\normalsize\bfseries}
  {}{0em}{}

\titlespacing*{\section}{0pt}{1.4em}{0.7em}
\titlespacing*{\subsection}{0pt}{1.1em}{0.45em}

% -----------------------------------------------------------------------------
% Links e corpo
% -----------------------------------------------------------------------------
\usepackage[hidelinks]{hyperref}
\hypersetup{
  colorlinks=true,
  linkcolor=corprincipal,
  urlcolor=corsecundaria,
  citecolor=corprincipal
}

\color{texto}
\setlength{\parindent}{0pt}
\setlength{\parskip}{0.4em}
\setstretch{1.12}
